\documentclass[12pt,a4paper]{report}

\usepackage[T1]{fontenc}
\usepackage[french]{babel}
\usepackage{lmodern}
\usepackage{geometry}
\usepackage{graphicx}
\usepackage{amsmath, amssymb}
\usepackage{hyperref}
\usepackage{fancyhdr}
\usepackage{setspace}
\usepackage{titlesec}
\usepackage{tocloft}

\geometry{margin=2.5cm}
\onehalfspacing

\pagestyle{fancy}
\fancyhf{}
\fancyhead[L]{TP 3G \& LTE}
\fancyhead[R]{\leftmark}
\fancyfoot[C]{\thepage}
\renewcommand{\headrulewidth}{0.4pt}
\setlength{\headheight}{14.5pt}

\title{
    \vspace{2cm}
    \textbf{\Huge TP 3G \& LTE} \\
    \vspace{0.5cm}
}
\author{
    Flotte Dylan \\
    Computer Science \\
    HES-SO Master
}
\date{\today}

\begin{document}

\maketitle
\thispagestyle{empty}
\newpage

\newpage

\tableofcontents
\newpage

\chapter{Code Python}

\section{config.py}
\begin{verbatim}
# Paramètres de connexion
HOST = "127.0.0.1"
PORT = 5000

# Clé secrète Ki (stockée dans USIM et AuC)
# En production, cette clé devrait être stockée de manière sécurisée
Ki = b"SuperSecretKey123"

# Paramètres de debug
DEBUG = True

\end{verbatim}

\section{crypto.py}
\begin{verbatim}
import os
import hashlib
from config import DEBUG


def generate_rand():
    """
    Génère un nombre aléatoire RAND de 16 octets
    Utilisé par le serveur pour le challenge d'authentification
    """
    rand = os.urandom(16)
    if DEBUG:
        print(f"[CRYPTO] RAND généré: {rand.hex()}")
    return rand


def compute_sres(rand, ki):
    """
    Calcule la réponse signée SRES à partir de RAND et Ki
    SRES = premiers 4 octets de SHA256(RAND || Ki)
    
    Args:
        rand: Le challenge RAND (16 octets)
        ki: La clé secrète Ki
    
    Returns:
        SRES (4 octets)
    """
    sres = hashlib.sha256(rand + ki).digest()[:4]
    if DEBUG:
        print(f"[CRYPTO] SRES calculé: {sres.hex()}")
    return sres


def derive_kc(rand, ki):
    """
    Dérive la clé de session Kc à partir de RAND et Ki
    Kc = SHA256(Ki || RAND)
    
    Args:
        rand: Le challenge RAND (16 octets)
        ki: La clé secrète Ki
    
    Returns:
        Kc (32 octets)
    """
    kc = hashlib.sha256(ki + rand).digest()
    if DEBUG:
        print(f"[CRYPTO] Clé de session Kc dérivée: {kc.hex()}")
    return kc


def xor_cipher(data, key):
    """
    Chiffrement/déchiffrement XOR avec clé
    
    Args:
        data: Données à chiffrer/déchiffrer
        key: Clé de chiffrement
    
    Returns:
        Données chiffrées/déchiffrées
    """
    result = bytes([data[i] ^ key[i % len(key)] for i in range(len(data))])
    if DEBUG:
        print(f"[CRYPTO] XOR cipher appliqué - Taille: {len(data)} octets")
    return result


def compute_auth_token(rand, ki):
    """
    Calcule le token d'authentification du serveur
    Token = SHA256(RAND || Ki)
    
    Args:
        rand: Le challenge RAND
        ki: La clé secrète Ki
    
    Returns:
        Token d'authentification (32 octets)
    """
    token = hashlib.sha256(rand + ki).digest()
    if DEBUG:
        print(f"[CRYPTO] Token d'authentification calculé: {token.hex()}")
    return token

\end{verbatim}

\section{session.py}
\begin{verbatim}
from config import Ki, DEBUG
from crypto import compute_sres, derive_kc, xor_cipher, compute_auth_token


class Session:
    """
    Classe gérant l'état d'une session d'authentification
    """
    
    def __init__(self):
        self.rand = None
        self.kc = None
        self.authenticated = False
        self.server_authenticated = False
        if DEBUG:
            print("[SESSION] Session initialisée")
    
    def set_rand(self, rand):
        """
        Enregistre le RAND reçu/généré
        """
        self.rand = rand
        if DEBUG:
            print(f"[SESSION] RAND enregistré: {rand.hex()}")
    
    def authenticate_client(self, received_sres):
        """
        Authentifie le client en vérifiant le SRES reçu
        
        Args:
            received_sres: SRES reçu du client
        
        Returns:
            True si authentifié, False sinon
        """
        if self.rand is None:
            if DEBUG:
                print("[SESSION] Erreur: RAND non défini")
            return False
        
        expected_sres = compute_sres(self.rand, Ki)
        self.authenticated = (received_sres == expected_sres)
        
        if DEBUG:
            if self.authenticated:
                print("[SESSION] Client authentifié avec succès")
            else:
                print("[SESSION] Échec d'authentification du client")
                print(f"[SESSION]   SRES attendu: {expected_sres.hex()}")
                print(f"[SESSION]   SRES reçu: {received_sres.hex()}")
        
        return self.authenticated
    
    def authenticate_server(self, received_token):
        """
        Authentifie le serveur en vérifiant le token reçu
        
        Args:
            received_token: Token reçu du serveur
        
        Returns:
            True si authentifié, False sinon
        """
        if self.rand is None:
            if DEBUG:
                print("[SESSION] Erreur: RAND non défini")
            return False
        
        expected_token = compute_auth_token(self.rand, Ki)
        self.server_authenticated = (received_token == expected_token)
        
        if DEBUG:
            if self.server_authenticated:
                print("[SESSION] Serveur authentifié avec succès")
            else:
                print("[SESSION] Échec d'authentification du serveur")
                print(f"[SESSION]   Token attendu: {expected_token.hex()}")
                print(f"[SESSION]   Token reçu: {received_token.hex()}")
        
        return self.server_authenticated
    
    def establish_session_key(self):
        """
        Établit la clé de session Kc
        
        Returns:
            La clé de session Kc
        """
        if self.rand is None:
            if DEBUG:
                print("[SESSION] Erreur: RAND non défini")
            return None
        
        self.kc = derive_kc(self.rand, Ki)
        if DEBUG:
            print(f"[SESSION] Clé de session établie: {self.kc.hex()}")
        return self.kc
    
    def encrypt_data(self, data):
        """
        Chiffre des données avec la clé de session
        
        Args:
            data: Données à chiffrer
        
        Returns:
            Données chiffrées
        """
        if self.kc is None:
            if DEBUG:
                print("[SESSION] Erreur: Clé de session non établie")
            return None
        
        if DEBUG:
            print(f"[SESSION] Chiffrement de {len(data)} octets")
        
        return xor_cipher(data, self.kc)
    
    def decrypt_data(self, cipher_data):
        """
        Déchiffre des données avec la clé de session
        
        Args:
            cipher_data: Données chiffrées
        
        Returns:
            Données déchiffrées
        """
        if self.kc is None:
            if DEBUG:
                print("[SESSION] Erreur: Clé de session non établie")
            return None
        
        if DEBUG:
            print(f"[SESSION] Déchiffrement de {len(cipher_data)} octets")
        
        return xor_cipher(cipher_data, self.kc)
    
    def get_auth_token(self):
        """
        Génère le token d'authentification du serveur
        
        Returns:
            Token d'authentification
        """
        if self.rand is None:
            if DEBUG:
                print("[SESSION] Erreur: RAND non défini")
            return None
        
        return compute_auth_token(self.rand, Ki)
    
    def is_authenticated(self):
        """
        Vérifie si la session est authentifiée
        
        Returns:
            True si authentifié, False sinon
        """
        return self.authenticated
    
    def is_server_authenticated(self):
        """
        Vérifie si le serveur est authentifié
        
        Returns:
            True si authentifié, False sinon
        """
        return self.server_authenticated

\end{verbatim}

\section{server.py}
\begin{verbatim}
import socket
from config import HOST, PORT, DEBUG
from crypto import generate_rand
from session import Session

def main():
    # Initialisation de la session
    session = Session()
    
    # Création du socket serveur
    server = socket.socket(socket.AF_INET, socket.SOCK_STREAM)
    server.bind((HOST, PORT))
    server.listen(1)
    
    print(f"[SERVEUR] En attente de connexion sur {HOST}:{PORT}...")
    print()
    
    conn, addr = server.accept()
    print(f"[SERVEUR] Client connecté depuis: {addr}")
    print()
    
    try:
        # ÉTAPE 1: Génération et envoi du RAND
        RAND = generate_rand()
        session.set_rand(RAND)
        conn.send(RAND)
        print(f"[SERVEUR] RAND envoyé au client: {RAND.hex()}")
        print()
        
        # ÉTAPE 2: Réception et vérification du SRES
        print("[SERVEUR] Attente de la réponse SRES du client...")
        client_sres = conn.recv(1024)
        print(f"[SERVEUR] SRES reçu: {client_sres.hex()}")
        
        if not session.authenticate_client(client_sres):
            print("[SERVEUR] ÉCHEC: Authentification du client échouée")
            conn.close()
            server.close()
            return
        
        print("[SERVEUR] Client authentifié avec succès")
        print()
        
        # ÉTAPE 3: Authentification du serveur auprès du client
        token = session.get_auth_token()
        conn.send(token)
        print(f"[SERVEUR] Token d'authentification envoyé: {token.hex()}")
        print()
        
        # ÉTAPE 4: Établissement de la clé de session
        Kc = session.establish_session_key()
        print(f"[SERVEUR] Clé de session Kc établie: {Kc.hex()}")
        print()
        
        # ÉTAPE 5: Réception des données chiffrées
        print("[SERVEUR] Attente des données du client...")
        cipher_data = conn.recv(1024)
        print(f"[SERVEUR] Données chiffrées reçues ({len(cipher_data)} octets): {cipher_data.hex()}")
        
        plain = session.decrypt_data(cipher_data)
        print(f"[SERVEUR] Données déchiffrées: {plain}")
        print()
        
        # ÉTAPE 6: Envoi de la réponse chiffrée
        message = b"Accuse de reception"
        print(f"[SERVEUR] Message à envoyer: {message}")
        cipher = session.encrypt_data(message)
        print(f"[SERVEUR] Message chiffré: {cipher.hex()}")
        conn.send(cipher)
        print("[SERVEUR] Réponse envoyée au client")
        print()
        
    except Exception as e:
        print(f"[SERVEUR] Erreur: {e}")
    finally:
        conn.close()
        server.close()
        print("[SERVEUR] Connexion fermée")

if __name__ == "__main__":
    main()
\end{verbatim}

\section{client.py}

\begin{verbatim}
import socket
from config import HOST, PORT, DEBUG
from crypto import compute_sres
from session import Session

def main():
    # Initialisation de la session
    session = Session()
    
    # Connexion au serveur
    client = socket.socket(socket.AF_INET, socket.SOCK_STREAM)
    
    print(f"[CLIENT] Connexion au serveur {HOST}:{PORT}...")
    
    try:
        client.connect((HOST, PORT))
        print("[CLIENT] Connecté au serveur")
        print()
        
        # ÉTAPE 1: Réception du RAND
        print("[CLIENT] Attente du RAND du serveur...")
        RAND = client.recv(1024)
        session.set_rand(RAND)
        print(f"[CLIENT] RAND reçu: {RAND.hex()}")
        print()
        
        # ÉTAPE 2: Calcul et envoi du SRES
        print("[CLIENT] Calcul du SRES...")
        from config import Ki
        SRES = compute_sres(RAND, Ki)
        print(f"[CLIENT] SRES calculé: {SRES.hex()}")
        client.send(SRES)
        print("[CLIENT] SRES envoyé au serveur")
        print()
        
        # ÉTAPE 3: Authentification du serveur
        print("[CLIENT] Attente du token d'authentification...")
        token = client.recv(1024)
        print(f"[CLIENT] Token reçu: {token.hex()}")
        
        if not session.authenticate_server(token):
            print("[CLIENT] ÉCHEC: Serveur non authentifié")
            client.close()
            return
        
        print("[CLIENT] Serveur authentifié avec succès")
        print()
        
        # ÉTAPE 4: Établissement de la clé de session
        Kc = session.establish_session_key()
        print(f"[CLIENT] Clé de session Kc établie: {Kc.hex()}")
        print()
        
        # ÉTAPE 5: Envoi de données chiffrées
        # Paquet de données (~256 bits)
        data = b"a3f7b2c1d9e845607f2c1a3b9d4e6f8a01c3b5d7e9f2a4c6b8d0e1f3a5c7b9d0"
        print(f"[CLIENT] Données à envoyer ({len(data)} octets): {data}")
        
        cipher = session.encrypt_data(data)
        print(f"[CLIENT] Données chiffrées: {cipher.hex()}")
        client.send(cipher)
        print("[CLIENT] Données envoyées au serveur")
        print()
        
        # ÉTAPE 6: Réception de la réponse
        print("[CLIENT] Attente de la réponse du serveur...")
        cipher_resp = client.recv(1024)
        print(f"[CLIENT] Réponse chiffrée reçue: {cipher_resp.hex()}")
        
        resp = session.decrypt_data(cipher_resp)
        print(f"[CLIENT] Réponse déchiffrée: {resp}")
        print()
        
    except Exception as e:
        print(f"[CLIENT] Erreur: {e}")
    finally:
        client.close()
        print("[CLIENT] Connexion fermée")

if __name__ == "__main__":
    main()
\end{verbatim}

\chapter{Capture du fonctionnement}

\begin{figure}[h]
	\centering
    \includegraphics[width=1\textwidth]{Cap1.png}
    \caption{Capture fonctionnement côté serveur}
    \label{fig:monimage}
\end{figure}

\begin{figure}[h]
	\centering
    \includegraphics[width=1\textwidth]{Cap2.png}
    \caption{Capture fonctionnement côté client}
    \label{fig:monimage}
\end{figure}

\end{document}
